\subsection{Suprema and Infima}

\begin{remark*}[Exposition]
Suprema and infima are sharp bounds.  A supremum is the least upper
bound; an infimum is the greatest lower bound.  The set remains required
to be nonempty throughout this section.
\end{remark*}

\begin{figure}[htbp]
  \centering
  \begin{tikzpicture}[
  x=1cm,
  y=1cm,
  >=Latex,
  axis/.style={thick},
  interval/.style={line width=2pt,teal!70!black},
  openpt/.style={circle, draw=teal!70!black, fill=white, thick, inner sep=2.2pt},
  bound/.style={circle, fill=blue!60!black, inner sep=2pt},
  lbl/.style={font=\small},
  note/.style={font=\footnotesize,text=gray!70!black}
]

\draw[axis] (-3.2,0) -- (3.2,0);
\draw[interval] (-1.4,0) -- (1.4,0);
\node[openpt] at (-1.4,0) {};
\node[openpt] at (1.4,0) {};

\node[bound] at (-1.4,0) {};
\node[bound] at (1.4,0) {};
\draw[blue!60!black,thick] (-1.4,-0.28) -- (-1.4,0.28);
\draw[blue!60!black,thick] (1.4,-0.28) -- (1.4,0.28);

\node[lbl,blue!60!black,above] at (-1.4,0.28) {$\inf A$};
\node[lbl,blue!60!black,above] at (1.4,0.28) {$\sup A$};
\node[lbl,teal!70!black,below] at (0,-0.28) {$A$};

\draw[gray!60,thick] (-2.6,-0.16) -- (-2.6,0.16);
\draw[gray!60,thick] (2.6,-0.16) -- (2.6,0.16);
\node[note,below] at (-2.6,-0.16) {lower bound};
\node[note,below] at (2.6,-0.16) {upper bound};

\draw[->,gray!70] (-1.4,0.65) -- (-2.45,0.2);
\draw[->,gray!70] (1.4,0.65) -- (2.45,0.2);
\node[note] at (0,0.9) {sharp bounds need not be attained};

\end{tikzpicture}

  \caption{Suprema and infima are sharp bounds, whether or not they are elements of the set.}
  \label{fig:suprema-infima-combinations}
\end{figure}

\begin{figure}[htbp]
  \centering
  \begin{tikzpicture}[
  >=Stealth,
  x=4.0cm,
  y=1.0cm,
  every node/.style={font=\small}
]
  \draw[->, gray!70] (-0.2,0) -- (1.25,0) node[right] {$\mathbb{R}$};
  \draw[line width=4pt, blue!35] (0,0) -- (1,0);
  \draw[blue, thick] (0,0) circle (2.2pt);
  \draw[blue, thick] (1,0) circle (2.2pt);
  \node[below] at (0,0) {$0=\inf(0,1)$};
  \node[below] at (1,0) {$1=\sup(0,1)$};
  \node[above] at (0.5,0.12) {$A=(0,1)$};
  \node[align=center] at (0.5,-0.9)
    {The sharp bounds sit at the endpoints,\\but neither endpoint belongs to $A$.};
\end{tikzpicture}

  \caption{For the open interval $(0,1)$, the supremum and infimum are
  endpoint bounds that are not attained by the set.}
  \label{fig:supremum-infimum-nonattainment}
\end{figure}

\begin{definitionbox}{Supremum}
\begin{definition}[Supremum]
\label{def:supremum}
Let $(S,\leq)$ be a partially ordered set, let $A\subseteq S$, and let
$s\in S$. The element $s$ is the \emph{supremum} of $A$, written
$s=\sup A$, if
\[
  (\forall a\in A)(a\leq s)
\]
and
\[
  (\forall u\in S)
  \left[
    (\forall a\in A)(a\leq u)
    \Longrightarrow
    s\leq u
  \right].
\]
Thus, $s$ is the least upper bound of $A$ in $S$.
\end{definition}
\end{definitionbox}

\begin{remark*}[Standard quantified statement]
\[
  A\neq\varnothing
  \quad\text{and}\quad
  (\forall a\in A)(a\leq s)
  \quad\text{and}\quad
  (\forall u\in\mathbb{R})\bigl[(\forall a\in A)(a\leq u)\Rightarrow s\leq u\bigr].
\]
\end{remark*}

\begin{remark*}[Predicate reading]
\[
  P=\mathsf{OrderedSet}(\mathbb{R},\leq),\qquad A\subseteq\mathbb{R},\qquad
  \operatorname{LeastUpperBound}(s,A,P)
  :=
  \operatorname{UpperBound}(s,A,P)
  \land
  (\forall u\in\mathbb{R})
  \bigl(\operatorname{UpperBound}(u,A,P)\Rightarrow s\leq u\bigr).
\]
\end{remark*}

\begin{remark*}[Negated quantified statement]
\[
  A=\varnothing
  \quad\text{or}\quad
  (\exists a\in A)(s<a)
  \quad\text{or}\quad
  (\exists u\in\mathbb{R})\bigl[(\forall a\in A)(a\leq u)\ \text{and}\ u<s\bigr].
\]
\end{remark*}

\begin{remark*}[Negation predicate reading]
\[
  P=\mathsf{OrderedSet}(\mathbb{R},\leq),\qquad A\subseteq\mathbb{R},\qquad
  \neg\operatorname{LeastUpperBound}(s,A,P)
  \Longleftrightarrow
  A=\varnothing
  \lor
  \neg\operatorname{UpperBound}(s,A,P)
  \lor
  (\exists u\in\mathbb{R})
  \bigl(\operatorname{UpperBound}(u,A,P)\land u<s\bigr).
\]
\end{remark*}

\begin{remark*}[Failure modes]
\begin{description}
  \item[Exposition.]
  The set may be empty; $s$ may fail to be an upper bound; or $s$ may be an
  upper bound that is not least among upper bounds.

  \item[\(\operatorname{LeastUpperBound} / \operatorname{UpperBound}.\)]
  The first clause is a domain failure.  The second is a bound failure.  The
  third is a sharpness failure witnessed by a strictly smaller upper bound.
  \[
    A=\varnothing
    \quad\text{or}\quad
    (\exists a\in A)(s<a)
    \quad\text{or}\quad
    (\exists u\in\mathbb{R})\bigl[(\forall a\in A)(a\leq u)\ \text{and}\ u<s\bigr].
  \]
  \[
    \neg\operatorname{LeastUpperBound}(s,A,P)
    \Longleftrightarrow
    A=\varnothing
    \lor
    \neg\operatorname{UpperBound}(s,A,P)
    \lor
    (\exists u\in\mathbb{R})
    \bigl(\operatorname{UpperBound}(u,A,P)\land u<s\bigr).
  \]
\end{description}
\end{remark*}

\begin{remark*}[Interpretation]
A supremum is the least upper bound: it combines upper-bound control with leastness among all upper bounds.
\end{remark*}

\begin{remark*}[Exposition]
Nonemptiness is a standing hypothesis for real suprema in this section.  The
empty set is bounded above in $\mathbb{R}$, but it has no real supremum:
every real number is an upper bound of $\varnothing$, and $\mathbb{R}$ has no
least element.
\end{remark*}

\begin{dependencies}
\begin{itemize}
  \item \hyperref[def:ordered-set]{Ordered Set}
  \item \hyperref[def:subset]{Subset}
  \item \hyperref[def:real-upper-bound]{Upper Bound}
\end{itemize}
\end{dependencies}

\begin{definitionbox}{Infimum}
\begin{definition}[Infimum]
\label{def:infimum}
Let $(S,\leq)$ be a partially ordered set, let $A\subseteq S$, and let
$i\in S$. The element $i$ is the \emph{infimum} of $A$, written
$i=\inf A$, if
\[
  (\forall a\in A)(i\leq a)
\]
and
\[
  (\forall \ell\in S)
  \left[
    (\forall a\in A)(\ell\leq a)
    \Longrightarrow
    \ell\leq i
  \right].
\]
Thus, $i$ is the greatest lower bound of $A$ in $S$.
\end{definition}
\end{definitionbox}

\begin{remark*}[Standard quantified statement]
\[
  A\neq\varnothing
  \quad\text{and}\quad
  (\forall a\in A)(i\leq a)
  \quad\text{and}\quad
  (\forall \ell\in\mathbb{R})\bigl[(\forall a\in A)(\ell\leq a)\Rightarrow \ell\leq i\bigr].
\]
\end{remark*}

\begin{remark*}[Predicate reading]
\[
  P=\mathsf{OrderedSet}(\mathbb{R},\leq),\qquad A\subseteq\mathbb{R},\qquad
  \operatorname{GreatestLowerBound}(i,A,P)
  :=
  \operatorname{LowerBound}(i,A,P)
  \land
  (\forall \ell\in\mathbb{R})
  \bigl(\operatorname{LowerBound}(\ell,A,P)\Rightarrow \ell\leq i\bigr).
\]
\end{remark*}

\begin{remark*}[Negated quantified statement]
\[
  A=\varnothing
  \quad\text{or}\quad
  (\exists a\in A)(a<i)
  \quad\text{or}\quad
  (\exists \ell\in\mathbb{R})\bigl[(\forall a\in A)(\ell\leq a)\ \text{and}\ i<\ell\bigr].
\]
\end{remark*}

\begin{remark*}[Negation predicate reading]
\[
  P=\mathsf{OrderedSet}(\mathbb{R},\leq),\qquad A\subseteq\mathbb{R},\qquad
  \neg\operatorname{GreatestLowerBound}(i,A,P)
  \Longleftrightarrow
  A=\varnothing
  \lor
  \neg\operatorname{LowerBound}(i,A,P)
  \lor
  (\exists \ell\in\mathbb{R})
  \bigl(\operatorname{LowerBound}(\ell,A,P)\land i<\ell\bigr).
\]
\end{remark*}

\begin{remark*}[Failure modes]
\begin{description}
  \item[Exposition.]
  The set may be empty; $i$ may fail to be a lower bound; or $i$ may be a
  lower bound that is not greatest among lower bounds.

  \item[\(\operatorname{GreatestLowerBound} / \operatorname{LowerBound}.\)]
  The first clause is a domain failure.  The second is a bound failure.  The
  third is a sharpness failure witnessed by a strictly larger lower bound.
  \[
    A=\varnothing
    \quad\text{or}\quad
    (\exists a\in A)(a<i)
    \quad\text{or}\quad
    (\exists \ell\in\mathbb{R})\bigl[(\forall a\in A)(\ell\leq a)\ \text{and}\ i<\ell\bigr].
  \]
  \[
    \neg\operatorname{GreatestLowerBound}(i,A,P)
    \Longleftrightarrow
    A=\varnothing
    \lor
    \neg\operatorname{LowerBound}(i,A,P)
    \lor
    (\exists \ell\in\mathbb{R})
    \bigl(\operatorname{LowerBound}(\ell,A,P)\land i<\ell\bigr).
  \]
\end{description}
\end{remark*}

\begin{remark*}[Interpretation]
An infimum is the greatest lower bound: it combines lower-bound control with greatestness among all lower bounds.
\end{remark*}

\begin{remark*}[Exposition]
Nonemptiness is a standing hypothesis for real infima in this section.  The
empty set is bounded below in $\mathbb{R}$, but it has no real infimum:
every real number is a lower bound of $\varnothing$, and $\mathbb{R}$ has no
greatest element.
\end{remark*}

\begin{dependencies}
\begin{itemize}
  \item \hyperref[def:ordered-set]{Ordered Set}
  \item \hyperref[def:subset]{Subset}
  \item \hyperref[def:real-lower-bound]{Lower Bound}
\end{itemize}
\end{dependencies}

\begin{propositionbox}{Uniqueness of the Supremum}
\begin{proposition}[Uniqueness of the Supremum]
\label{prop:uniqueness-of-the-supremum}
Let $(S,\leq)$ be a partially ordered set and let $A\subseteq S$. If
$s,t\in S$ are both suprema of $A$, then
\[
  s=t.
\]

\hyperref[prf:uniqueness-of-the-supremum]{Go to proof.}
\end{proposition}
\end{propositionbox}

\begin{remark*}[Standard quantified statement]
\[
  A\neq\varnothing
  \quad\text{and}\quad
  \operatorname{LeastUpperBound}(s,A,P)
  \quad\text{and}\quad
  \operatorname{LeastUpperBound}(t,A,P)
  \quad\Longrightarrow\quad
  s=t.
\]
\end{remark*}

\begin{remark*}[Predicate reading]
\[
  S\in\mathbf{Set},\qquad A\subseteq S,\qquad P=\mathsf{OrderedSet}(S,\leq),\qquad
  \operatorname{LeastUpperBound}(s,A,P)\land\operatorname{LeastUpperBound}(t,A,P)
  \Longrightarrow s=t.
\]
\end{remark*}

\begin{remark*}[Negated quantified statement]
\[
  A\neq\varnothing
  \quad\text{and}\quad
  \operatorname{LeastUpperBound}(s,A,P)
  \quad\text{and}\quad
  \operatorname{LeastUpperBound}(t,A,P)
  \quad\text{and}\quad
  s\neq t.
\]
\end{remark*}

\begin{remark*}[Negation predicate reading]
\[
  S\in\mathbf{Set},\qquad A\subseteq S,\qquad P=\mathsf{OrderedSet}(S,\leq),\qquad
  \operatorname{LeastUpperBound}(s,A,P)\land
  \operatorname{LeastUpperBound}(t,A,P)\land s\neq t.
\]
\end{remark*}

\begin{remark*}[Failure modes]
\begin{description}
  \item[Exposition.]
  Failure would require two distinct least upper bounds for the same
  nonempty set.

  \item[\(\operatorname{LeastUpperBound}.\)]
  Each candidate is an upper bound and each lies below every upper bound;
  applying these minimality clauses to the other candidate gives both
  $s\leq t$ and $t\leq s$.
  \[
    A\neq\varnothing
    \quad\text{and}\quad
    \operatorname{LeastUpperBound}(s,A,P)
    \quad\text{and}\quad
    \operatorname{LeastUpperBound}(t,A,P)
    \quad\text{and}\quad
    s\neq t.
  \]
  \[
    P=\mathsf{OrderedSet}(S,\leq),\qquad
    \operatorname{LeastUpperBound}(s,A,P)\land
    \operatorname{LeastUpperBound}(t,A,P)\land s\neq t.
  \]
\end{description}
\end{remark*}

\begin{remark*}[Interpretation]
The statement records how the supremum predicate controls uniqueness or upper-bound behavior in the ambient order.
\end{remark*}

\begin{dependencies}
\begin{itemize}
  \item \hyperref[def:ordered-set]{Ordered Set}
  \item \hyperref[def:subset]{Subset}
  \item \hyperref[def:supremum]{Supremum}
\end{itemize}
\end{dependencies}

\begin{propositionbox}{Uniqueness of the Infimum}
\begin{proposition}[Uniqueness of the Infimum]
\label{prop:uniqueness-of-the-infimum}
Let $(S,\leq)$ be a partially ordered set and let $A\subseteq S$. If
$i,j\in S$ are both infima of $A$, then
\[
  i=j.
\]

\hyperref[prf:uniqueness-of-the-infimum]{Go to proof.}
\end{proposition}
\end{propositionbox}

\begin{remark*}[Standard quantified statement]
\[
  A\neq\varnothing
  \quad\text{and}\quad
  \operatorname{GreatestLowerBound}(i,A,P)
  \quad\text{and}\quad
  \operatorname{GreatestLowerBound}(j,A,P)
  \quad\Longrightarrow\quad
  i=j.
\]
\end{remark*}

\begin{remark*}[Predicate reading]
\[
  S\in\mathbf{Set},\qquad A\subseteq S,\qquad P=\mathsf{OrderedSet}(S,\leq),\qquad
  \operatorname{GreatestLowerBound}(i,A,P)\land\operatorname{GreatestLowerBound}(j,A,P)
  \Longrightarrow i=j.
\]
\end{remark*}

\begin{remark*}[Negated quantified statement]
\[
  A\neq\varnothing
  \quad\text{and}\quad
  \operatorname{GreatestLowerBound}(i,A,P)
  \quad\text{and}\quad
  \operatorname{GreatestLowerBound}(j,A,P)
  \quad\text{and}\quad
  i\neq j.
\]
\end{remark*}

\begin{remark*}[Negation predicate reading]
\[
  S\in\mathbf{Set},\qquad A\subseteq S,\qquad P=\mathsf{OrderedSet}(S,\leq),\qquad
  \operatorname{GreatestLowerBound}(i,A,P)\land
  \operatorname{GreatestLowerBound}(j,A,P)\land i\neq j.
\]
\end{remark*}

\begin{remark*}[Failure modes]
\begin{description}
  \item[Exposition.]
  Failure would require two distinct greatest lower bounds for the same
  nonempty set.

  \item[\(\operatorname{GreatestLowerBound}.\)]
  Each candidate is a lower bound and each lies above every lower bound;
  applying these maximality clauses to the other candidate gives both
  $i\leq j$ and $j\leq i$.
  \[
    A\neq\varnothing
    \quad\text{and}\quad
    \operatorname{GreatestLowerBound}(i,A,P)
    \quad\text{and}\quad
    \operatorname{GreatestLowerBound}(j,A,P)
    \quad\text{and}\quad
    i\neq j.
  \]
  \[
    P=\mathsf{OrderedSet}(S,\leq),\qquad
    \operatorname{GreatestLowerBound}(i,A,P)\land
    \operatorname{GreatestLowerBound}(j,A,P)\land i\neq j.
  \]
\end{description}
\end{remark*}

\begin{remark*}[Interpretation]
The statement records how the infimum predicate controls uniqueness or lower-bound behavior in the ambient order.
\end{remark*}

\begin{dependencies}
\begin{itemize}
  \item \hyperref[def:ordered-set]{Ordered Set}
  \item \hyperref[def:subset]{Subset}
  \item \hyperref[def:infimum]{Infimum}
\end{itemize}
\end{dependencies}

\begin{propositionbox}{Upper Bound Property of the Supremum}
\begin{proposition}[Upper-Bound Property of the Supremum]
\label{prop:upper-bound-property-of-supremum}
Let $(S,\leq)$ be a partially ordered set, let $A\subseteq S$, and suppose
that $s=\sup A$ exists in $S$. Then $s$ is an upper bound of $A$; that is,
\[
  (\forall x\in A)(x\leq s).
\]

\hyperref[prf:upper-bound-property-of-supremum]{Go to proof.}
\end{proposition}
\end{propositionbox}

\begin{remark*}[Standard quantified statement]
\[
  A\neq\varnothing,
  \quad (\forall x\in A)(x\leq s),
  \quad\text{and}\quad
  (\forall u\in S)\bigl[(\forall x\in A)(x\leq u)\Rightarrow s\leq u\bigr]
  \quad\Longrightarrow\quad
  (\forall x\in A)(x\leq s).
\]
\end{remark*}

\begin{remark*}[Predicate reading]
\[
  S\in\mathbf{Set},\qquad A\subseteq S,\qquad P=\mathsf{OrderedSet}(S,\leq),\qquad
  \operatorname{LeastUpperBound}(s,A,P)\Longrightarrow
  \operatorname{UpperBound}(s,A,P).
\]
\end{remark*}

\begin{remark*}[Negated quantified statement]
\[
  \operatorname{LeastUpperBound}(s,A,P)
  \quad\text{and}\quad
  (\exists x\in A)(s<x).
\]
\end{remark*}

\begin{remark*}[Negation predicate reading]
\[
  \operatorname{LeastUpperBound}(s,A,P)
  \quad\text{and}\quad
  (\exists x\in A)(s<x).
\]

The negation predicate reading is the predicate-level form of the displayed negated statement.
\end{remark*}

\begin{remark*}[Failure modes]
\begin{description}
  \item[Exposition.]
  Failure would mean that a supremum is not an upper bound.

  
\[
\operatorname{LeastUpperBound}(s,A,P)
  \quad\text{and}\quad
  (\exists x\in A)(s<x).
\]
  \item[\(\operatorname{LeastUpperBound}.\)]
  This is exactly the first clause in the definition of supremum.
  \[
    \operatorname{LeastUpperBound}(s,A,P)
    \quad\text{and}\quad
    (\exists x\in A)(s<x).
  \]

\[
    \operatorname{LeastUpperBound}(s,A,P)
    \quad\text{and}\quad
    (\exists x\in A)(s<x).
  \]
\end{description}
\end{remark*}

\begin{remark*}[Interpretation]
The statement records how the supremum predicate controls uniqueness or upper-bound behavior in the ambient order.
\end{remark*}

\begin{dependencies}
\begin{itemize}
  \item \hyperref[def:ordered-set]{Ordered Set}
  \item \hyperref[def:subset]{Subset}
  \item \hyperref[def:supremum]{Supremum}
\end{itemize}
\end{dependencies}

\begin{propositionbox}{Lower Bound Property of the Infimum}
\begin{proposition}[Lower-Bound Property of the Infimum]
\label{prop:lower-bound-property-of-infimum}
Let $(S,\leq)$ be a partially ordered set, let $A\subseteq S$, and suppose
that $i=\inf A$ exists in $S$. Then $i$ is a lower bound of $A$; that is,
\[
  (\forall x\in A)(i\leq x).
\]

\hyperref[prf:lower-bound-property-of-infimum]{Go to proof.}
\end{proposition}
\end{propositionbox}

\begin{remark*}[Standard quantified statement]
\[
  A\neq\varnothing,
  \quad (\forall x\in A)(i\leq x),
  \quad\text{and}\quad
  (\forall \ell\in S)\bigl[(\forall x\in A)(\ell\leq x)\Rightarrow \ell\leq i\bigr]
  \quad\Longrightarrow\quad
  (\forall x\in A)(i\leq x).
\]
\end{remark*}

\begin{remark*}[Predicate reading]
\[
  S\in\mathbf{Set},\qquad A\subseteq S,\qquad P=\mathsf{OrderedSet}(S,\leq),\qquad
  \operatorname{GreatestLowerBound}(i,A,P)\Longrightarrow
  \operatorname{LowerBound}(i,A,P).
\]
\end{remark*}

\begin{remark*}[Negated quantified statement]
\[
  \operatorname{GreatestLowerBound}(i,A,P)
  \quad\text{and}\quad
  (\exists x\in A)(x<i).
\]
\end{remark*}

\begin{remark*}[Negation predicate reading]
\[
  \operatorname{GreatestLowerBound}(i,A,P)
  \quad\text{and}\quad
  (\exists x\in A)(x<i).
\]

The negation predicate reading is the predicate-level form of the displayed negated statement.
\end{remark*}

\begin{remark*}[Failure modes]
\begin{description}
  \item[Exposition.]
  Failure would mean that an infimum is not a lower bound.

  
\[
\operatorname{GreatestLowerBound}(i,A,P)
  \quad\text{and}\quad
  (\exists x\in A)(x<i).
\]
  \item[\(\operatorname{GreatestLowerBound}.\)]
  This is exactly the first clause in the definition of infimum.
  \[
    \operatorname{GreatestLowerBound}(i,A,P)
    \quad\text{and}\quad
    (\exists x\in A)(x<i).
  \]

\[
    \operatorname{GreatestLowerBound}(i,A,P)
    \quad\text{and}\quad
    (\exists x\in A)(x<i).
  \]
\end{description}
\end{remark*}

\begin{remark*}[Interpretation]
The statement records how the infimum predicate controls uniqueness or lower-bound behavior in the ambient order.
\end{remark*}

\begin{dependencies}
\begin{itemize}
  \item \hyperref[def:ordered-set]{Ordered Set}
  \item \hyperref[def:subset]{Subset}
  \item \hyperref[def:infimum]{Infimum}
\end{itemize}
\end{dependencies}

\begin{lemmabox}{Subset Inclusion Preserves Upper Bounds}
\begin{lemma}[Subset Inclusion Preserves Upper Bounds]
\label{lem:subset-inclusion-preserves-upper-bounds}
Let $(S,\leq)$ be a partially ordered set, let $A,B\subseteq S$, and suppose
that $A\subseteq B$. If $u\in S$ is an upper bound of $B$, then $u$ is an
upper bound of $A$; equivalently,
\[
  \left[(\forall b\in B)(b\leq u)\right]
  \Longrightarrow
  \left[(\forall a\in A)(a\leq u)\right].
\]

\hyperref[prf:subset-inclusion-preserves-upper-bounds]{Go to proof.}
\end{lemma}
\end{lemmabox}

\begin{remark*}[Standard quantified statement]
\[
  A\subseteq B,\quad
  (\forall b\in B)(b\leq u)
  \quad\Longrightarrow\quad
  (\forall a\in A)(a\leq u).
\]
\end{remark*}

\begin{remark*}[Predicate reading]
\[
  S\in\mathbf{Set},\qquad A\subseteq B\subseteq S,\qquad P=\mathsf{OrderedSet}(S,\leq),
\]
\[
  \operatorname{UpperBound}(u,B,P)
  \Longrightarrow \operatorname{UpperBound}(u,A,P).
\]
\end{remark*}

\begin{remark*}[Negated quantified statement]
\[
  A\subseteq B,\quad \operatorname{UpperBound}(u,B,P),
  \quad\text{and}\quad
  (\exists x\in A)(u<x).
\]
\end{remark*}

\begin{remark*}[Negation predicate reading]
\[
  A\subseteq B,\quad \operatorname{UpperBound}(u,B,P),
  \quad\text{and}\quad
  (\exists x\in A)(u<x).
\]

The negation predicate reading is the predicate-level form of the displayed negated statement.
\end{remark*}

\begin{remark*}[Failure modes]
\begin{description}
  \item[Exposition.]
  Failure would require a counterexample in the smaller set even though
  the larger set is already bounded by $u$.

  
\[
A\subseteq B,\quad \operatorname{UpperBound}(u,B,P),
  \quad\text{and}\quad
  (\exists x\in A)(u<x).
\]
  \item[\(\operatorname{UpperBound}.\)]
  Any counterexample $x\in A$ is also an element of $B$, contradicting the
  upper-bound hypothesis for $B$.
  \[
    A\subseteq B,\quad \operatorname{UpperBound}(u,B,P),
    \quad\text{and}\quad
    (\exists x\in A)(u<x).
  \]

\[
    A\subseteq B,\quad \operatorname{UpperBound}(u,B,P),
    \quad\text{and}\quad
    (\exists x\in A)(u<x).
  \]
\end{description}
\end{remark*}

\begin{remark*}[Interpretation]
A supremum is the least upper bound: it combines upper-bound control with leastness among all upper bounds.
\end{remark*}

\begin{dependencies}
\begin{itemize}
  \item \hyperref[def:ordered-set]{Ordered Set}
  \item \hyperref[def:subset]{Subset}
  \item \hyperref[def:real-upper-bound]{Upper Bound}
\end{itemize}
\end{dependencies}

\begin{lemmabox}{Subset Inclusion Preserves Lower Bounds}
\begin{lemma}[Subset Inclusion Preserves Lower Bounds]
\label{lem:subset-inclusion-preserves-lower-bounds}
Let $(S,\leq)$ be a partially ordered set, let $A,B\subseteq S$, and suppose
that $A\subseteq B$. If $\ell\in S$ is a lower bound of $B$, then $\ell$ is
a lower bound of $A$; equivalently,
\[
  \left[(\forall b\in B)(\ell\leq b)\right]
  \Longrightarrow
  \left[(\forall a\in A)(\ell\leq a)\right].
\]

\hyperref[prf:subset-inclusion-preserves-lower-bounds]{Go to proof.}
\end{lemma}
\end{lemmabox}

\begin{remark*}[Standard quantified statement]
\[
  A\subseteq B,\quad
  (\forall b\in B)(\ell\leq b)
  \quad\Longrightarrow\quad
  (\forall a\in A)(\ell\leq a).
\]
\end{remark*}

\begin{remark*}[Predicate reading]
\[
  S\in\mathbf{Set},\qquad A\subseteq B\subseteq S,\qquad P=\mathsf{OrderedSet}(S,\leq),
\]
\[
  \operatorname{LowerBound}(\ell,B,P)
  \Longrightarrow \operatorname{LowerBound}(\ell,A,P).
\]
\end{remark*}

\begin{remark*}[Negated quantified statement]
\[
  A\subseteq B,\quad \operatorname{LowerBound}(\ell,B,P),
  \quad\text{and}\quad
  (\exists x\in A)(x<\ell).
\]
\end{remark*}

\begin{remark*}[Negation predicate reading]
\[
  A\subseteq B,\quad \operatorname{LowerBound}(\ell,B,P),
  \quad\text{and}\quad
  (\exists x\in A)(x<\ell).
\]

The negation predicate reading is the predicate-level form of the displayed negated statement.
\end{remark*}

\begin{remark*}[Failure modes]
\begin{description}
  \item[Exposition.]
  Failure would require a smaller-set counterexample to a lower bound that
  already works on the larger set.

  
\[
A\subseteq B,\quad \operatorname{LowerBound}(\ell,B,P),
  \quad\text{and}\quad
  (\exists x\in A)(x<\ell).
\]
  \item[\(\operatorname{LowerBound}.\)]
  Any counterexample $x\in A$ lies in $B$, contradicting the lower-bound
  hypothesis for $B$.
  \[
    A\subseteq B,\quad \operatorname{LowerBound}(\ell,B,P),
    \quad\text{and}\quad
    (\exists x\in A)(x<\ell).
  \]

\[
    A\subseteq B,\quad \operatorname{LowerBound}(\ell,B,P),
    \quad\text{and}\quad
    (\exists x\in A)(x<\ell).
  \]
\end{description}
\end{remark*}

\begin{remark*}[Interpretation]
An infimum is the greatest lower bound: it combines lower-bound control with greatestness among all lower bounds.
\end{remark*}

\begin{dependencies}
\begin{itemize}
  \item \hyperref[def:ordered-set]{Ordered Set}
  \item \hyperref[def:subset]{Subset}
  \item \hyperref[def:real-lower-bound]{Lower Bound}
\end{itemize}
\end{dependencies}

\begin{theorembox}{Supremum Is Monotone Under Inclusion}
\begin{theorem}[Supremum Is Monotone Under Inclusion]
\label{thm:supremum-is-monotone-under-inclusion}
Let $A,B\subseteq\mathbb{R}$ be nonempty and bounded above. If
$A\subseteq B$, then
\[
  \sup A\leq\sup B.
\]

\hyperref[prf:supremum-is-monotone-under-inclusion]{Go to proof.}
\end{theorem}
\end{theorembox}

\begin{remark*}[Standard quantified statement]
\[
  A\neq\varnothing,\quad B\neq\varnothing,\quad
  A\subseteq B,\quad
  \operatorname{LeastUpperBound}(s_A,A,P),\quad
  \operatorname{LeastUpperBound}(s_B,B,P)
  \quad\Longrightarrow\quad s_A\leq s_B.
\]
\end{remark*}

\begin{remark*}[Predicate reading]
\[
  P=\mathsf{OrderedSet}(\mathbb{R},\leq),\qquad A\subseteq B\subseteq\mathbb{R},\qquad
  A\subseteq B\land\operatorname{LeastUpperBound}(s_A,A,P)\land\operatorname{LeastUpperBound}(s_B,B,P)
  \Longrightarrow s_A\leq s_B.
\]
\end{remark*}

\begin{remark*}[Negated quantified statement]
\[
  A\subseteq B,\quad \operatorname{LeastUpperBound}(s_A,A,P),\quad
  \operatorname{LeastUpperBound}(s_B,B,P),\quad s_B<s_A.
\]
\end{remark*}

\begin{remark*}[Negation predicate reading]
\[
  A\subseteq B,\quad \operatorname{LeastUpperBound}(s_A,A,P),\quad
  \operatorname{LeastUpperBound}(s_B,B,P),\quad s_B<s_A.
\]

The negation predicate reading is the predicate-level form of the displayed negated statement.
\end{remark*}

\begin{remark*}[Failure modes]
\begin{description}
  \item[Exposition.]
  Failure would mean enlarging a nonempty set lowers its least upper bound.

  
\[
A\subseteq B,\quad \operatorname{LeastUpperBound}(s_A,A,P),\quad
  \operatorname{LeastUpperBound}(s_B,B,P),\quad s_B<s_A.
\]
  \item[\(\operatorname{LeastUpperBound}.\)]
  The value $\sup B$ bounds $B$, hence bounds $A$; since $\sup A$ is least
  among upper bounds of $A$, $\sup A\leq\sup B$.
  \[
    A\subseteq B,\quad \operatorname{LeastUpperBound}(s_A,A,P),\quad
    \operatorname{LeastUpperBound}(s_B,B,P),\quad s_B<s_A.
  \]

\[
    A\subseteq B,\quad \operatorname{LeastUpperBound}(s_A,A,P),\quad
    \operatorname{LeastUpperBound}(s_B,B,P),\quad s_B<s_A.
  \]
\end{description}
\end{remark*}

\begin{remark*}[Interpretation]
The statement records how the supremum predicate controls uniqueness or upper-bound behavior in the ambient order.
\end{remark*}

\begin{dependencies}
\begin{itemize}
  \item \hyperref[def:reals]{Real Numbers}
  \item \hyperref[def:ordered-set]{Ordered Set}
  \item \hyperref[def:subset]{Subset}
  \item \hyperref[def:supremum]{Supremum}
  \item \hyperref[lem:subset-inclusion-preserves-upper-bounds]{Subset Inclusion Preserves Upper Bounds}
\end{itemize}
\end{dependencies}

\begin{theorembox}{Infimum Is Monotone Under Inclusion}
\begin{theorem}[Infimum Is Antitone Under Inclusion]
\label{thm:infimum-is-monotone-under-inclusion}
Let $A,B\subseteq\mathbb{R}$ be nonempty and bounded below. If
$A\subseteq B$, then
\[
  \inf B\leq\inf A.
\]

\hyperref[prf:infimum-is-monotone-under-inclusion]{Go to proof.}
\end{theorem}
\end{theorembox}

\begin{remark*}[Standard quantified statement]
\[
  A\neq\varnothing,\quad B\neq\varnothing,\quad
  A\subseteq B,\quad
  \operatorname{GreatestLowerBound}(i_A,A,P),\quad
  \operatorname{GreatestLowerBound}(i_B,B,P)
  \quad\Longrightarrow\quad i_B\leq i_A.
\]
\end{remark*}

\begin{remark*}[Predicate reading]
\[
  P=\mathsf{OrderedSet}(\mathbb{R},\leq),\qquad A\subseteq B\subseteq\mathbb{R},\qquad
  A\subseteq B\land\operatorname{GreatestLowerBound}(i_A,A,P)\land\operatorname{GreatestLowerBound}(i_B,B,P)
  \Longrightarrow i_B\leq i_A.
\]
\end{remark*}

\begin{remark*}[Negated quantified statement]
\[
  A\subseteq B,\quad \operatorname{GreatestLowerBound}(i_A,A,P),\quad
  \operatorname{GreatestLowerBound}(i_B,B,P),\quad i_A<i_B.
\]
\end{remark*}

\begin{remark*}[Negation predicate reading]
\[
  A\subseteq B,\quad \operatorname{GreatestLowerBound}(i_A,A,P),\quad
  \operatorname{GreatestLowerBound}(i_B,B,P),\quad i_A<i_B.
\]

The negation predicate reading is the predicate-level form of the displayed negated statement.
\end{remark*}

\begin{remark*}[Failure modes]
\begin{description}
  \item[Exposition.]
  Failure would mean enlarging a nonempty set raises its greatest lower
  bound.

  
\[
A\subseteq B,\quad \operatorname{GreatestLowerBound}(i_A,A,P),\quad
  \operatorname{GreatestLowerBound}(i_B,B,P),\quad i_A<i_B.
\]
  \item[\(\operatorname{GreatestLowerBound}.\)]
  The value $\inf B$ bounds $B$ below, hence bounds $A$ below; since
  $\inf A$ is greatest among lower bounds of $A$, $\inf B\leq\inf A$.
  \[
    A\subseteq B,\quad \operatorname{GreatestLowerBound}(i_A,A,P),\quad
    \operatorname{GreatestLowerBound}(i_B,B,P),\quad i_A<i_B.
  \]

\[
    A\subseteq B,\quad \operatorname{GreatestLowerBound}(i_A,A,P),\quad
    \operatorname{GreatestLowerBound}(i_B,B,P),\quad i_A<i_B.
  \]
\end{description}
\end{remark*}

\begin{remark*}[Interpretation]
The statement records how the infimum predicate controls uniqueness or lower-bound behavior in the ambient order.
\end{remark*}

\begin{dependencies}
\begin{itemize}
  \item \hyperref[def:reals]{Real Numbers}
  \item \hyperref[def:ordered-set]{Ordered Set}
  \item \hyperref[def:subset]{Subset}
  \item \hyperref[def:infimum]{Infimum}
  \item \hyperref[lem:subset-inclusion-preserves-lower-bounds]{Subset Inclusion Preserves Lower Bounds}
\end{itemize}
\end{dependencies}

\begin{propositionbox}{Upper Bound Comparison with the Supremum}
\begin{proposition}[Upper-Bound Comparison with the Supremum]
\label{prop:upper-bound-comparison-with-the-supremum}
Let $A\subseteq\mathbb{R}$ be nonempty and bounded above, and let
$s=\sup A$. For every $u\in\mathbb{R}$,
\[
  u\text{ is an upper bound of }A
  \Longleftrightarrow
  s\leq u.
\]
Equivalently,
\[
  (\forall a\in A)(a\leq u)
  \Longleftrightarrow
  s\leq u.
\]

\hyperref[prf:upper-bound-comparison-with-the-supremum]{Go to proof.}
\end{proposition}
\end{propositionbox}

\begin{remark*}[Standard quantified statement]
\[
  A\neq\varnothing,\quad \operatorname{LeastUpperBound}(s,A,P)
  \quad\Longrightarrow\quad
  (\forall u\in\mathbb{R})
  \bigl[\operatorname{UpperBound}(u,A,P)\Longleftrightarrow s\leq u\bigr].
\]
\end{remark*}

\begin{remark*}[Predicate reading]
\[
  P=\mathsf{OrderedSet}(\mathbb{R},\leq),\qquad A\subseteq\mathbb{R},\qquad
  \operatorname{LeastUpperBound}(s,A,P)\Longrightarrow
  \forall u\,(\operatorname{UpperBound}(u,A,P)\leftrightarrow s\leq u).
\]
\end{remark*}

\begin{remark*}[Negated quantified statement]
\[
  \operatorname{LeastUpperBound}(s,A,P)
  \quad\text{and}\quad
  (\exists u)\bigl[(\operatorname{UpperBound}(u,A,P)\land u<s)
  \lor(s\leq u\land\neg\operatorname{UpperBound}(u,A,P))\bigr].
\]
\end{remark*}

\begin{remark*}[Negation predicate reading]
\[
  \operatorname{LeastUpperBound}(s,A,P)
  \quad\text{and}\quad
  (\exists u)\bigl[(\operatorname{UpperBound}(u,A,P)\land u<s)
  \lor(s\leq u\land\neg\operatorname{UpperBound}(u,A,P))\bigr].
\]

The negation predicate reading is the predicate-level form of the displayed negated statement.
\end{remark*}

\begin{remark*}[Failure modes]
\begin{description}
  \item[Exposition.]
  Failure would mean an upper bound below the supremum, or a number above
  the supremum that fails to bound the set.

  
\[
\operatorname{LeastUpperBound}(s,A,P)
  \quad\text{and}\quad
  (\exists u)\bigl[(\operatorname{UpperBound}(u,A,P)\land u<s)
  \lor(s\leq u\land\neg\operatorname{UpperBound}(u,A,P))\bigr].
\]
  \item[\(\operatorname{LeastUpperBound} / \operatorname{UpperBound}.\)]
  Leastness gives the forward implication.  The fact that $s$ bounds $A$
  gives the reverse implication by transitivity.
  \[
    \operatorname{LeastUpperBound}(s,A,P)
    \quad\text{and}\quad
    (\exists u)\bigl[(\operatorname{UpperBound}(u,A,P)\land u<s)
    \lor(s\leq u\land\neg\operatorname{UpperBound}(u,A,P))\bigr].
  \]

\[
    \operatorname{LeastUpperBound}(s,A,P)
    \quad\text{and}\quad
    (\exists u)\bigl[(\operatorname{UpperBound}(u,A,P)\land u<s)
    \lor(s\leq u\land\neg\operatorname{UpperBound}(u,A,P))\bigr].
  \]
\end{description}
\end{remark*}

\begin{remark*}[Interpretation]
A supremum is the least upper bound: it combines upper-bound control with leastness among all upper bounds.
\end{remark*}

\begin{dependencies}
\begin{itemize}
  \item \hyperref[def:reals]{Real Numbers}
  \item \hyperref[def:ordered-set]{Ordered Set}
  \item \hyperref[def:subset]{Subset}
  \item \hyperref[def:real-upper-bound]{Upper Bound}
  \item \hyperref[def:supremum]{Supremum}
\end{itemize}
\end{dependencies}

\begin{propositionbox}{Lower Bound Comparison with the Infimum}
\begin{proposition}[Lower-Bound Comparison with the Infimum]
\label{prop:lower-bound-comparison-with-the-infimum}
Let $A\subseteq\mathbb{R}$ be nonempty and bounded below, and let
$i=\inf A$. For every $\ell\in\mathbb{R}$,
\[
  \ell\text{ is a lower bound of }A
  \Longleftrightarrow
  \ell\leq i.
\]
Equivalently,
\[
  (\forall a\in A)(\ell\leq a)
  \Longleftrightarrow
  \ell\leq i.
\]

\hyperref[prf:lower-bound-comparison-with-the-infimum]{Go to proof.}
\end{proposition}
\end{propositionbox}

\begin{remark*}[Standard quantified statement]
\[
  A\neq\varnothing,\quad \operatorname{GreatestLowerBound}(i,A,P)
  \quad\Longrightarrow\quad
  (\forall \ell\in\mathbb{R})
  \bigl[\operatorname{LowerBound}(\ell,A,P)\Longleftrightarrow \ell\leq i\bigr].
\]
\end{remark*}

\begin{remark*}[Predicate reading]
\[
  P=\mathsf{OrderedSet}(\mathbb{R},\leq),\qquad A\subseteq\mathbb{R},\qquad
  \operatorname{GreatestLowerBound}(i,A,P)\Longrightarrow
  \forall \ell\,(\operatorname{LowerBound}(\ell,A,P)\leftrightarrow \ell\leq i).
\]
\end{remark*}

\begin{remark*}[Negated quantified statement]
\[
  \operatorname{GreatestLowerBound}(i,A,P)
  \quad\text{and}\quad
  (\exists \ell)\bigl[(\operatorname{LowerBound}(\ell,A,P)\land i<\ell)
  \lor(\ell\leq i\land\neg\operatorname{LowerBound}(\ell,A,P))\bigr].
\]
\end{remark*}

\begin{remark*}[Negation predicate reading]
\[
  \operatorname{GreatestLowerBound}(i,A,P)
  \quad\text{and}\quad
  (\exists \ell)\bigl[(\operatorname{LowerBound}(\ell,A,P)\land i<\ell)
  \lor(\ell\leq i\land\neg\operatorname{LowerBound}(\ell,A,P))\bigr].
\]

The negation predicate reading is the predicate-level form of the displayed negated statement.
\end{remark*}

\begin{remark*}[Failure modes]
\begin{description}
  \item[Exposition.]
  Failure would mean a lower bound above the infimum, or a number below the
  infimum that fails to bound the set below.

  
\[
\operatorname{GreatestLowerBound}(i,A,P)
  \quad\text{and}\quad
  (\exists \ell)\bigl[(\operatorname{LowerBound}(\ell,A,P)\land i<\ell)
  \lor(\ell\leq i\land\neg\operatorname{LowerBound}(\ell,A,P))\bigr].
\]
  \item[\(\operatorname{GreatestLowerBound} / \operatorname{LowerBound}.\)]
  Greatestness gives the forward implication.  The fact that $i$ bounds
  $A$ below gives the reverse implication by transitivity.
  \[
    \operatorname{GreatestLowerBound}(i,A,P)
    \quad\text{and}\quad
    (\exists \ell)\bigl[(\operatorname{LowerBound}(\ell,A,P)\land i<\ell)
    \lor(\ell\leq i\land\neg\operatorname{LowerBound}(\ell,A,P))\bigr].
  \]

\[
    \operatorname{GreatestLowerBound}(i,A,P)
    \quad\text{and}\quad
    (\exists \ell)\bigl[(\operatorname{LowerBound}(\ell,A,P)\land i<\ell)
    \lor(\ell\leq i\land\neg\operatorname{LowerBound}(\ell,A,P))\bigr].
  \]
\end{description}
\end{remark*}

\begin{remark*}[Interpretation]
An infimum is the greatest lower bound: it combines lower-bound control with greatestness among all lower bounds.
\end{remark*}

\begin{dependencies}
\begin{itemize}
  \item \hyperref[def:reals]{Real Numbers}
  \item \hyperref[def:ordered-set]{Ordered Set}
  \item \hyperref[def:subset]{Subset}
  \item \hyperref[def:real-lower-bound]{Lower Bound}
  \item \hyperref[def:infimum]{Infimum}
\end{itemize}
\end{dependencies}

\begin{propositionbox}{Every Element Lies Below the Supremum}
\begin{proposition}[Every Element Lies Below the Supremum]
\label{prop:every-element-lies-below-the-supremum}
Let $A\subseteq\mathbb{R}$ be nonempty and bounded above, and let
$s=\sup A$. Then
\[
  (\forall x\in A)(x\leq s).
\]

\hyperref[prf:every-element-lies-below-the-supremum]{Go to proof.}
\end{proposition}
\end{propositionbox}

\begin{remark*}[Standard quantified statement]
\[
  A\neq\varnothing,\quad \operatorname{LeastUpperBound}(s,A,P)
  \quad\Longrightarrow\quad
  (\forall x\in A)(x\leq s).
\]
\end{remark*}

\begin{remark*}[Predicate reading]
\[
  P=\mathsf{OrderedSet}(\mathbb{R},\leq),\qquad A\subseteq\mathbb{R},\qquad
  \operatorname{LeastUpperBound}(s,A,P)\Longrightarrow \operatorname{UpperBound}(s,A,P).
\]
\end{remark*}

\begin{remark*}[Negated quantified statement]
\[
  \operatorname{LeastUpperBound}(s,A,P)\quad\text{and}\quad(\exists x\in A)(s<x).
\]
\end{remark*}

\begin{remark*}[Negation predicate reading]
\[
  \operatorname{LeastUpperBound}(s,A,P)\quad\text{and}\quad(\exists x\in A)(s<x).
\]

The negation predicate reading is the predicate-level form of the displayed negated statement.
\end{remark*}

\begin{remark*}[Failure modes]
\begin{description}
  \item[Exposition.]
  Failure would mean the supremum is not an upper bound.

  
\[
  \operatorname{LeastUpperBound}(s,A,P)\quad\text{and}\quad(\exists x\in A)(s<x).
\]
  \item[\(\operatorname{LeastUpperBound}.\)]
  This is the upper-bound clause in the definition of supremum.
  \[
    \operatorname{LeastUpperBound}(s,A,P)\quad\text{and}\quad(\exists x\in A)(s<x).
  \]

\[
    \operatorname{LeastUpperBound}(s,A,P)\quad\text{and}\quad(\exists x\in A)(s<x).
  \]
\end{description}
\end{remark*}

\begin{remark*}[Interpretation]
The statement records how the supremum predicate controls uniqueness or upper-bound behavior in the ambient order.
\end{remark*}

\begin{dependencies}
\begin{itemize}
  \item \hyperref[def:reals]{Real Numbers}
  \item \hyperref[def:ordered-set]{Ordered Set}
  \item \hyperref[def:subset]{Subset}
  \item \hyperref[def:supremum]{Supremum}
\end{itemize}
\end{dependencies}

\begin{propositionbox}{Every Element Lies Above the Infimum}
\begin{proposition}[Every Element Lies Above the Infimum]
\label{prop:every-element-lies-above-the-infimum}
Let $A\subseteq\mathbb{R}$ be nonempty and bounded below, and let
$i=\inf A$. Then
\[
  (\forall x\in A)(i\leq x).
\]

\hyperref[prf:every-element-lies-above-the-infimum]{Go to proof.}
\end{proposition}
\end{propositionbox}

\begin{remark*}[Standard quantified statement]
\[
  A\neq\varnothing,\quad \operatorname{GreatestLowerBound}(i,A,P)
  \quad\Longrightarrow\quad
  (\forall x\in A)(i\leq x).
\]
\end{remark*}

\begin{remark*}[Predicate reading]
\[
  P=\mathsf{OrderedSet}(\mathbb{R},\leq),\qquad A\subseteq\mathbb{R},\qquad
  \operatorname{GreatestLowerBound}(i,A,P)\Longrightarrow \operatorname{LowerBound}(i,A,P).
\]
\end{remark*}

\begin{remark*}[Negated quantified statement]
\[
  \operatorname{GreatestLowerBound}(i,A,P)\quad\text{and}\quad(\exists x\in A)(x<i).
\]
\end{remark*}

\begin{remark*}[Negation predicate reading]
\[
  \operatorname{GreatestLowerBound}(i,A,P)\quad\text{and}\quad(\exists x\in A)(x<i).
\]

The negation predicate reading is the predicate-level form of the displayed negated statement.
\end{remark*}

\begin{remark*}[Failure modes]
\begin{description}
  \item[Exposition.]
  Failure would mean the infimum is not a lower bound.

  
\[
  \operatorname{GreatestLowerBound}(i,A,P)\quad\text{and}\quad(\exists x\in A)(x<i).
\]
  \item[\(\operatorname{GreatestLowerBound}.\)]
  This is the lower-bound clause in the definition of infimum.
  \[
    \operatorname{GreatestLowerBound}(i,A,P)\quad\text{and}\quad(\exists x\in A)(x<i).
  \]

\[
    \operatorname{GreatestLowerBound}(i,A,P)\quad\text{and}\quad(\exists x\in A)(x<i).
  \]
\end{description}
\end{remark*}

\begin{remark*}[Interpretation]
The statement records how the infimum predicate controls uniqueness or lower-bound behavior in the ambient order.
\end{remark*}

\begin{dependencies}
\begin{itemize}
  \item \hyperref[def:reals]{Real Numbers}
  \item \hyperref[def:ordered-set]{Ordered Set}
  \item \hyperref[def:subset]{Subset}
  \item \hyperref[def:infimum]{Infimum}
\end{itemize}
\end{dependencies}

\begin{theorembox}{Infimum Is Less Than or Equal to the Supremum}
\begin{theorem}[Infimum Is Less Than or Equal to the Supremum]
\label{thm:infimum-less-than-supremum}
Let $A\subseteq\mathbb{R}$ be nonempty and bounded. Then
\[
  \inf A\leq\sup A.
\]

\hyperref[prf:infimum-less-than-supremum]{Go to proof.}
\end{theorem}
\end{theorembox}

\begin{remark*}[Standard quantified statement]
\[
  A\neq\varnothing,\quad
  \operatorname{GreatestLowerBound}(i,A,P),\quad
  \operatorname{LeastUpperBound}(s,A,P)
  \quad\Longrightarrow\quad
  i\leq s.
\]
\end{remark*}

\begin{remark*}[Predicate reading]
\[
  P=\mathsf{OrderedSet}(\mathbb{R},\leq),\qquad A\subseteq\mathbb{R},\qquad
  \operatorname{GreatestLowerBound}(i,A,P)\land\operatorname{LeastUpperBound}(s,A,P)
  \Longrightarrow i\leq s.
\]
\end{remark*}

\begin{remark*}[Negated quantified statement]
\[
  \operatorname{GreatestLowerBound}(i,A,P),\quad \operatorname{LeastUpperBound}(s,A,P),
  \quad\text{and}\quad s<i.
\]
\end{remark*}

\begin{remark*}[Negation predicate reading]
\[
  \operatorname{GreatestLowerBound}(i,A,P),\quad \operatorname{LeastUpperBound}(s,A,P),
  \quad\text{and}\quad s<i.
\]

The negation predicate reading is the predicate-level form of the displayed negated statement.
\end{remark*}

\begin{remark*}[Failure modes]
\begin{description}
  \item[Exposition.]
  Failure would put the greatest lower bound strictly above the least upper
  bound of the same nonempty set.

  
\[
\operatorname{GreatestLowerBound}(i,A,P),\quad \operatorname{LeastUpperBound}(s,A,P),
  \quad\text{and}\quad s<i.
\]
  \item[\(\operatorname{GreatestLowerBound} / \operatorname{LeastUpperBound}.\)]
  Choose $a\in A$.  Then $i\leq a$ because $i$ is a lower bound, and
  $a\leq s$ because $s$ is an upper bound.
  \[
    \operatorname{GreatestLowerBound}(i,A,P),\quad \operatorname{LeastUpperBound}(s,A,P),
    \quad\text{and}\quad s<i.
  \]

\[
    \operatorname{GreatestLowerBound}(i,A,P),\quad \operatorname{LeastUpperBound}(s,A,P),
    \quad\text{and}\quad s<i.
  \]
\end{description}
\end{remark*}

\begin{remark*}[Interpretation]
The statement records how the supremum predicate controls uniqueness or upper-bound behavior in the ambient order.
\end{remark*}

\begin{dependencies}
\begin{itemize}
  \item \hyperref[def:reals]{Real Numbers}
  \item \hyperref[def:ordered-set]{Ordered Set}
  \item \hyperref[def:subset]{Subset}
  \item \hyperref[def:infimum]{Infimum}
  \item \hyperref[def:supremum]{Supremum}
\end{itemize}
\end{dependencies}

\begin{propositionbox}{Proposition (Supremum Need Not Belong to the Set)}
\begin{proposition}[Supremum Need Not Belong to the Set]
\label{prop:supremum-need-not-belong-to-the-set}
There exists a nonempty bounded-above set $A\subseteq\mathbb{R}$ such that
\[
  \sup A\notin A.
\]

\hyperref[prf:supremum-need-not-belong-to-the-set]{Go to proof.}
\end{proposition}
\end{propositionbox}

\begin{remark*}[Standard quantified statement]
\[
  (\exists A\subseteq\mathbb{R})(\exists s\in\mathbb{R})
  \bigl[A\neq\varnothing\land \operatorname{LeastUpperBound}(s,A,P)\land s\notin A\bigr].
\]
\end{remark*}

\begin{remark*}[Predicate reading]
\[
  P=\mathsf{OrderedSet}(\mathbb{R},\leq),\qquad
  (\exists A\subseteq\mathbb{R})(\exists s)\bigl(\operatorname{LeastUpperBound}(s,A,P)\land s\notin A\bigr).
\]
\end{remark*}

\begin{remark*}[Interpretation]
A supremum is a sharp upper bound, but sharpness does not require the bound
to be attained by the set.
\end{remark*}

\begin{dependencies}
\begin{itemize}
  \item \hyperref[def:reals]{Real Numbers}
  \item \hyperref[def:ordered-set]{Ordered Set}
  \item \hyperref[def:subset]{Subset}
  \item \hyperref[def:supremum]{Supremum}
  \item \hyperref[thm:lub-property-implies-existence-of-suprema]{Least Upper Bound Property Implies Existence of Suprema}
\end{itemize}
\end{dependencies}

\begin{propositionbox}{Proposition (Infimum Need Not Belong to the Set)}
\begin{proposition}[Infimum Need Not Belong to the Set]
\label{prop:infimum-need-not-belong-to-the-set}
There exists a nonempty bounded-below set $A\subseteq\mathbb{R}$ such that
\[
  \inf A\notin A.
\]

\hyperref[prf:infimum-need-not-belong-to-the-set]{Go to proof.}
\end{proposition}
\end{propositionbox}

\begin{remark*}[Standard quantified statement]
\[
  (\exists A\subseteq\mathbb{R})(\exists i\in\mathbb{R})
  \bigl[A\neq\varnothing\land \operatorname{GreatestLowerBound}(i,A,P)\land i\notin A\bigr].
\]
\end{remark*}

\begin{remark*}[Predicate reading]
\[
  P=\mathsf{OrderedSet}(\mathbb{R},\leq),\qquad
  (\exists A\subseteq\mathbb{R})(\exists i)\bigl(\operatorname{GreatestLowerBound}(i,A,P)\land i\notin A\bigr).
\]
\end{remark*}

\begin{remark*}[Interpretation]
An infimum is a sharp lower bound, but sharpness does not require the bound
to be attained by the set.
\end{remark*}

\begin{dependencies}
\begin{itemize}
  \item \hyperref[def:reals]{Real Numbers}
  \item \hyperref[def:ordered-set]{Ordered Set}
  \item \hyperref[def:subset]{Subset}
  \item \hyperref[def:infimum]{Infimum}
  \item \hyperref[thm:glb-property-implies-existence-of-infima]{Greatest Lower Bound Property Implies Existence of Infima}
\end{itemize}
\end{dependencies}

\begin{propositionbox}{Order Comparison of Suprema}
\begin{proposition}[Order Comparison of Suprema]
\label{prop:order-comparison-of-suprema}
Let $A,B\subseteq\mathbb{R}$ be nonempty and bounded above. Suppose that
\[
  (\forall a\in A)(\exists b\in B)(a\leq b).
\]
Then
\[
  \sup A\leq\sup B.
\]

\hyperref[prf:order-comparison-of-suprema]{Go to proof.}
\end{proposition}
\end{propositionbox}

\begin{remark*}[Standard quantified statement]
\[
  A\neq\varnothing,\quad B\neq\varnothing,\quad
  \operatorname{LeastUpperBound}(s_A,A,P),\quad
  \operatorname{LeastUpperBound}(s_B,B,P),\quad
  (\forall a\in A)(\exists b\in B)(a\leq b)
  \quad\Longrightarrow\quad
  s_A\leq s_B.
\]
\end{remark*}

\begin{remark*}[Predicate reading]
\[
  P=\mathsf{OrderedSet}(\mathbb{R},\leq),\qquad
  A,B\subseteq\mathbb{R},\qquad
  [\forall x\in A\,\exists y\in B\,(x\leq y)]
  \land\operatorname{LeastUpperBound}(s_A,A,P)\land\operatorname{LeastUpperBound}(s_B,B,P)
  \Longrightarrow s_A\leq s_B.
\]
\end{remark*}

\begin{remark*}[Negated quantified statement]
\[
  [\forall x\in A\,\exists y\in B\,(x\leq y)]
  \quad\text{and}\quad
  \sup B<\sup A.
\]
\end{remark*}

\begin{remark*}[Negation predicate reading]
\[
  [\forall x\in A\,\exists y\in B\,(x\leq y)]
  \quad\text{and}\quad
  \sup B<\sup A.
\]

The negation predicate reading is the predicate-level form of the displayed negated statement.
\end{remark*}

\begin{remark*}[Failure modes]
\begin{description}
  \item[Exposition.]
  Failure would mean $A$ reaches higher in supremum than $B$ even though
  each element of $A$ is dominated by an element of $B$.

  
\[
[\forall x\in A\,\exists y\in B\,(x\leq y)]
  \quad\text{and}\quad
  \sup B<\sup A.
\]
  \item[\(\operatorname{LeastUpperBound}.\)]
  Every element of $B$ lies below $\sup B$, so every element of $A$ also
  lies below $\sup B$.  Thus $\sup B$ is an upper bound for $A$.
  \[
    [\forall x\in A\,\exists y\in B\,(x\leq y)]
    \quad\text{and}\quad
    \sup B<\sup A.
  \]

\[
    [\forall x\in A\,\exists y\in B\,(x\leq y)]
    \quad\text{and}\quad
    \sup B<\sup A.
  \]
\end{description}
\end{remark*}

\begin{remark*}[Interpretation]
The statement records how the supremum predicate controls uniqueness or upper-bound behavior in the ambient order.
\end{remark*}

\begin{dependencies}
\begin{itemize}
  \item \hyperref[def:reals]{Real Numbers}
  \item \hyperref[def:ordered-set]{Ordered Set}
  \item \hyperref[def:subset]{Subset}
  \item \hyperref[def:supremum]{Supremum}
  \item \hyperref[prop:every-element-lies-below-the-supremum]{Every Element Lies Below the Supremum}
\end{itemize}
\end{dependencies}

\begin{theorembox}{Least Upper Bound Property Implies Existence of Suprema}
\begin{theorem}[Least-Upper-Bound Property Implies Existence of Suprema]
\label{thm:lub-property-implies-existence-of-suprema}
Let $A\subseteq\mathbb{R}$ be nonempty and bounded above. Then there
exists a unique $s\in\mathbb{R}$ such that
\[
  s=\sup A.
\]
Equivalently,
\[
  (\forall A\subseteq\mathbb{R})
  \left[
    A\neq\varnothing
    \land
    (\exists u\in\mathbb{R})(\forall a\in A)(a\leq u)
    \Longrightarrow
    (\exists!s\in\mathbb{R})(s=\sup A)
  \right].
\]

\hyperref[prf:lub-property-implies-existence-of-suprema]{Go to proof.}
\end{theorem}
\end{theorembox}

\begin{remark*}[Standard quantified statement]
\[
  S\neq\varnothing
  \quad\text{and}\quad
  (\exists u\in\mathbb{R})(\forall x\in S)(x\leq u)
  \quad\Longrightarrow\quad
  (\exists s\in\mathbb{R})\operatorname{LeastUpperBound}(s,S,P).
\]
\end{remark*}

\begin{remark*}[Predicate reading]
\[
  P=\mathsf{OrderedSet}(\mathbb{R},\leq),\qquad S\subseteq\mathbb{R},\qquad
  S\neq\varnothing\land\operatorname{BoundedAbove}(S,P)
  \Longrightarrow
  (\exists s\in\mathbb{R})\operatorname{LeastUpperBound}(s,S,P).
\]
\end{remark*}

\begin{remark*}[Negated quantified statement]
\[
  S\neq\varnothing,\quad
  \operatorname{BoundedAbove}(S,P),
  \quad\text{and}\quad
  (\forall s\in\mathbb{R})\neg\operatorname{LeastUpperBound}(s,S,P).
\]
\end{remark*}

\begin{remark*}[Negation predicate reading]
\[
  S\neq\varnothing,\quad
  \operatorname{BoundedAbove}(S,P),
  \quad\text{and}\quad
  (\forall s\in\mathbb{R})\neg\operatorname{LeastUpperBound}(s,S,P).
\]

The negation predicate reading is the predicate-level form of the displayed negated statement.
\end{remark*}

\begin{remark*}[Failure modes]
\begin{description}
  \item[Exposition.]
  In $\mathbb{R}$ the theorem cannot fail once the set is nonempty and
  bounded above; outside complete ordered settings, the existence clause
  can fail.

  
\[
S\neq\varnothing,\quad
  \operatorname{BoundedAbove}(S,P),
  \quad\text{and}\quad
  (\forall s\in\mathbb{R})\neg\operatorname{LeastUpperBound}(s,S,P).
\]
  \item[\(\operatorname{BoundedAbove} / \operatorname{LeastUpperBound}.\)]
  The hypotheses separate domain nonemptiness from one-sided boundedness.
  The conclusion asserts existence of the sharp upper bound.
  \[
    S\neq\varnothing,\quad
    \operatorname{BoundedAbove}(S,P),
    \quad\text{and}\quad
    (\forall s\in\mathbb{R})\neg\operatorname{LeastUpperBound}(s,S,P).
  \]

\[
    S\neq\varnothing,\quad
    \operatorname{BoundedAbove}(S,P),
    \quad\text{and}\quad
    (\forall s\in\mathbb{R})\neg\operatorname{LeastUpperBound}(s,S,P).
  \]
\end{description}
\end{remark*}

\begin{remark*}[Interpretation]
A supremum is the least upper bound: it combines upper-bound control with leastness among all upper bounds.
\end{remark*}

\begin{dependencies}
\begin{itemize}
  \item \hyperref[def:reals]{Real Numbers}
  \item \hyperref[def:ordered-set]{Ordered Set}
  \item \hyperref[def:subset]{Subset}
  \item \hyperref[def:least-upper-bound-property]{Least Upper Bound Property}
  \item \hyperref[def:real-upper-bound]{Upper Bound}
  \item \hyperref[def:bounded-above]{Bounded Above}
  \item \hyperref[def:supremum]{Supremum}
\end{itemize}
\end{dependencies}

\begin{theorembox}{Greatest Lower Bound Property Implies Existence of Infima}
\begin{theorem}[Greatest-Lower-Bound Property Implies Existence of Infima]
\label{thm:glb-property-implies-existence-of-infima}
Let $A\subseteq\mathbb{R}$ be nonempty and bounded below. Then there
exists a unique $i\in\mathbb{R}$ such that
\[
  i=\inf A.
\]
Equivalently,
\[
  (\forall A\subseteq\mathbb{R})
  \left[
    A\neq\varnothing
    \land
    (\exists \ell\in\mathbb{R})(\forall a\in A)(\ell\leq a)
    \Longrightarrow
    (\exists!i\in\mathbb{R})(i=\inf A)
  \right].
\]

\hyperref[prf:glb-property-implies-existence-of-infima]{Go to proof.}
\end{theorem}
\end{theorembox}

\begin{remark*}[Standard quantified statement]
\[
  S\neq\varnothing
  \quad\text{and}\quad
  (\exists \ell\in\mathbb{R})(\forall x\in S)(\ell\leq x)
  \quad\Longrightarrow\quad
  (\exists i\in\mathbb{R})\operatorname{GreatestLowerBound}(i,S,P).
\]
\end{remark*}

\begin{remark*}[Predicate reading]
\[
  P=\mathsf{OrderedSet}(\mathbb{R},\leq),\qquad S\subseteq\mathbb{R},\qquad
  S\neq\varnothing\land\operatorname{BoundedBelow}(S,P)
  \Longrightarrow
  (\exists i\in\mathbb{R})\operatorname{GreatestLowerBound}(i,S,P).
\]
\end{remark*}

\begin{remark*}[Negated quantified statement]
\[
  S\neq\varnothing,\quad
  \operatorname{BoundedBelow}(S,P),
  \quad\text{and}\quad
  (\forall i\in\mathbb{R})\neg\operatorname{GreatestLowerBound}(i,S,P).
\]
\end{remark*}

\begin{remark*}[Negation predicate reading]
\[
  S\neq\varnothing,\quad
  \operatorname{BoundedBelow}(S,P),
  \quad\text{and}\quad
  (\forall i\in\mathbb{R})\neg\operatorname{GreatestLowerBound}(i,S,P).
\]

The negation predicate reading is the predicate-level form of the displayed negated statement.
\end{remark*}

\begin{remark*}[Failure modes]
\begin{description}
  \item[Exposition.]
  In $\mathbb{R}$ the theorem cannot fail once the set is nonempty and
  bounded below.

  
\[
S\neq\varnothing,\quad
  \operatorname{BoundedBelow}(S,P),
  \quad\text{and}\quad
  (\forall i\in\mathbb{R})\neg\operatorname{GreatestLowerBound}(i,S,P).
\]
  \item[\(\operatorname{BoundedBelow} / \operatorname{GreatestLowerBound}.\)]
  Apply the least upper bound property to $-S$ or argue by the order-dual
  form of completeness.
  \[
    S\neq\varnothing,\quad
    \operatorname{BoundedBelow}(S,P),
    \quad\text{and}\quad
    (\forall i\in\mathbb{R})\neg\operatorname{GreatestLowerBound}(i,S,P).
  \]

\[
    S\neq\varnothing,\quad
    \operatorname{BoundedBelow}(S,P),
    \quad\text{and}\quad
    (\forall i\in\mathbb{R})\neg\operatorname{GreatestLowerBound}(i,S,P).
  \]
\end{description}
\end{remark*}

\begin{remark*}[Interpretation]
An infimum is the greatest lower bound: it combines lower-bound control with greatestness among all lower bounds.
\end{remark*}

\begin{dependencies}
\begin{itemize}
  \item \hyperref[def:reals]{Real Numbers}
  \item \hyperref[def:ordered-set]{Ordered Set}
  \item \hyperref[def:subset]{Subset}
  \item \hyperref[def:least-upper-bound-property]{Least Upper Bound Property}
  \item \hyperref[def:real-lower-bound]{Lower Bound}
  \item \hyperref[def:bounded-below]{Bounded Below}
  \item \hyperref[def:infimum]{Infimum}
\end{itemize}
\end{dependencies}

\begin{corollarybox}{Bounded Set Has Both a Supremum and an Infimum}
\begin{corollary}[Bounded Set Has Both a Supremum and an Infimum]
\label{cor:bounded-set-has-supremum-and-infimum}
Let $A\subseteq\mathbb{R}$ be nonempty and bounded. Then there exist unique
$s,i\in\mathbb{R}$ such that
\[
  s=\sup A \qquad\text{and}\qquad i=\inf A.
\]
\hyperref[prf:bounded-set-has-supremum-and-infimum]{\textit{Go to proof.}}
\end{corollary}
\end{corollarybox}

\begin{remark*}[Standard quantified statement]
\[
  A\neq\varnothing,\quad
  (\exists u\in\mathbb{R})(\forall a\in A)(a\leq u),\quad
  (\exists \ell\in\mathbb{R})(\forall a\in A)(\ell\leq a)
  \quad\Longrightarrow\quad
  (\exists! s\in\mathbb{R})\operatorname{LeastUpperBound}(s,A,P)
  \ \text{and}\
  (\exists! i\in\mathbb{R})\operatorname{GreatestLowerBound}(i,A,P).
\]
\end{remark*}

\begin{remark*}[Predicate reading]
\[
  P=\mathsf{OrderedSet}(\mathbb{R},\leq),\qquad A\subseteq\mathbb{R},\qquad
  A\neq\varnothing\land\operatorname{Bounded}(A,P)
  \Longrightarrow
  (\exists! s)\operatorname{LeastUpperBound}(s,A,P)\land(\exists! i)\operatorname{GreatestLowerBound}(i,A,P).
\]
\end{remark*}

\begin{remark*}[Negated quantified statement]
\[
  A\neq\varnothing,\quad \operatorname{Bounded}(A,P),
  \quad\text{and}\quad
  \bigl[(\forall s\in\mathbb{R})\neg\operatorname{LeastUpperBound}(s,A,P)
  \ \text{or}\
  (\forall i\in\mathbb{R})\neg\operatorname{GreatestLowerBound}(i,A,P)\bigr].
\]
\end{remark*}

\begin{remark*}[Negation predicate reading]
\[
  A\neq\varnothing,\quad \operatorname{Bounded}(A,P),
  \quad\text{and}\quad
  \bigl[(\forall s\in\mathbb{R})\neg\operatorname{LeastUpperBound}(s,A,P)
  \ \text{or}\
  (\forall i\in\mathbb{R})\neg\operatorname{GreatestLowerBound}(i,A,P)\bigr].
\]

The negation predicate reading is the predicate-level form of the displayed negated statement.
\end{remark*}

\begin{remark*}[Failure modes]
\begin{description}
  \item[Exposition.]
\[
  \mathbb{R}
\]
  In $\mathbb{R}$, neither existence clause can fail once the set is
  nonempty and bounded on the relevant side; this corollary simply records
  that both hold simultaneously when the set is bounded on both sides.

  
\[
A\neq\varnothing,\quad \operatorname{Bounded}(A,P),
  \quad\text{and}\quad
  \bigl[(\forall s\in\mathbb{R})\neg\operatorname{LeastUpperBound}(s,A,P)
  \ \text{or}\
  (\forall i\in\mathbb{R})\neg\operatorname{GreatestLowerBound}(i,A,P)\bigr].
\]
  \item[\(\operatorname{LeastUpperBound} / \operatorname{GreatestLowerBound}.\)]
\[
A\neq\varnothing,\quad \operatorname{Bounded}(A,P),
  \quad\text{and}\quad
  \bigl[(\forall s\in\mathbb{R})\neg\operatorname{LeastUpperBound}(s,A,P)
  \ \text{or}\
  (\forall i\in\mathbb{R})\neg\operatorname{GreatestLowerBound}(i,A,P)\bigr].
\]

\[
A\neq\varnothing,\quad \operatorname{Bounded}(A,P),
  \quad\text{and}\quad
  \bigl[(\forall s\in\mathbb{R})\neg\operatorname{LeastUpperBound}(s,A,P)
  \ \text{or}\
  (\forall i\in\mathbb{R})\neg\operatorname{GreatestLowerBound}(i,A,P)\bigr].
\]
  Apply \hyperref[thm:lub-property-implies-existence-of-suprema]{Least Upper
  Bound Property Implies Existence of Suprema} to the upper-bound half of
  $\operatorname{Bounded}(A,P)$, and
  \hyperref[thm:glb-property-implies-existence-of-infima]{Greatest Lower
  Bound Property Implies Existence of Infima} to the lower-bound half.
\end{description}
\end{remark*}

\begin{remark*}[Interpretation]
This is nothing more than the two existence theorems applied side by side:
a set that is bounded on both sides is exactly a set for which both the
supremum and the infimum are guaranteed to exist, each uniquely. It is
stated separately here because ``bounded'' (both-sided) is the hypothesis
that shows up constantly in later chapters, and it is worth having the
combined conclusion on hand as its own citable result rather than
re-deriving it from the two one-sided theorems every time.
\end{remark*}

\begin{dependencies}
\begin{itemize}
  \item \hyperref[def:bounded]{Bounded Set}
  \item \hyperref[thm:lub-property-implies-existence-of-suprema]{Least Upper Bound Property Implies Existence of Suprema}
  \item \hyperref[thm:glb-property-implies-existence-of-infima]{Greatest Lower Bound Property Implies Existence of Infima}
\end{itemize}
\end{dependencies}
